\scanpage{225}
\[
(9)\qquad \sum D_i\omega'_i=\overline\omega'\qquad\text{dans }\underline\Phi
\]
et la condition (7) signifie que l'on a, localement au dessus de $S$, une relation
\[
(10)\qquad \overline\omega\left(=\sum\overline D_{\deleted{\unsure{s}}i}\omega_i\right)=d\varphi,\qquad \varphi\in\Gamma(X,\mathcal O_X).
\]
Pour voir que cette relation implique bien la nullité de (3), il suffit de prouver
\[
(11)\qquad \sum_i D_{is}\int_{Z_s}\omega_{is}=\int_{Z_s}\overline\omega
\]
qui implique, si $\overline\omega=d\varphi$, que le premier mb est égal à $\int_{\partial Z_s}\varphi=0$ car $\partial Z_s=0$.

En fait, (11) provient de la relation générale
\[
(12)\qquad D_s\int_{Z_s}\omega_s=\int_{Z_s}\overline D_s\omega_s
\]
valable pour tt op. diff. hol. $D$ sur $S$, et tte forme \inserted{holomorphe fermée} $\omega_{\deleted{\unsure{i}}}$ sur $X$ [le fait que $\omega$ est fermée implique que $\int_{Z_s}\omega$ est fonction holomorphe en $s$, et donne un sens au premier mb de (12)]. Pour prouver cette
\margin{224}
